% !TEX program = lualatex
\documentclass[11pt,a4paper]{ltjsarticle}
\usepackage{icat-common}

% 定理環境の定義
\newtheorem{theorem}{定理}[section]
\newtheorem{definition}{定義}[section]
\newtheorem{axiom}{公理}[section]
\newtheorem{corollary}{系}[section]

% 独自マクロ
\renewcommand{\proofname}{\textbf{Proof.}}

\title{空数学における中道不動点コンビネータ $\Ymw$ の定式化と\\無明の幾何学的・証明論的実態}
\author{千葉将臣}
\date{2026-04-22}

\begin{document}
\maketitle

\begin{abstract}
本稿は、空数学 (Sunyata-Mathematics) の直観主義四句分別において要請される「自己参照の閉じ損ね」を、中道不動点コンビネータ $\Ymw$ として代数的に定式化するものである。本コンビネータが界論における「無明 (Avidy\=a)」 の発生ジェネレータとして機能し、いかにして無限のフラクタル構造（色）を展開するのかを論証する。さらに、その再帰的展開が第四段階のトポロジカルな散逸を経て、究極の真理保存則である自明な空コンビネータ $\Svoid$ へと相転移する機序、そこからの代数的な再起動、およびアルゴリズム情報理論 (AIT) に基づく解脱とカオスの決定不能性を証明する。
\end{abstract}

\section{序言}
直観主義四句分別において、真理は静的な状態ではなく構成プロセスとして現れる。本体系を駆動するためには、知覚の最小単位がフラクタル則を満たす「無明（自己参照の閉じ損ね）」として存在することが前提とされてきた。本稿では、この動的な発生プロセスをコンビネータ論理へと拡張し、メタ否定演算子 $\metanegation$ と中道等価 $\mweq$ によって構成される二つの極限的コンビネータ（発生器 $\Ymw$ と吸収器 $\Svoid$）を定義することで、世界の顕現と解脱のライフサイクルを完全に代数化する。

\section{定義}
\begin{definition}[メタ否定演算子と中道等価]
対象に対して境界を生成し、次元を非対称に引き上げるメタ観測を $\metanegation$ と定義する。また、厳密な自己同一性ではなく、フラクタル収束の軌道が一致することを以て等価とみなす非実体的な関係演算子を中道等価 $\mweq$ と定義する。
\end{definition}

\begin{definition}[中道不動点コンビネータ $\Ymw$]
メタ否定演算子 $\metanegation$ に対し、以下の中道等価を満たすコンビネータを $\Ymw$ と定義する。
\begin{equation}
\Ymw \metanegation \mweq \metanegation (\Ymw \metanegation)
\end{equation}
これは古典的ラムダ計算におけるYコンビネータの中道的拡張であり、界論における「無明ジェネレータ」として機能する。
\end{definition}

\begin{definition}[空コンビネータ $\Svoid$]
任意の関数（推論の横棒） $f$ を適用しても、常に自己同一性を返す自明なコンビネータを $\Svoid$ と定義する。
\begin{equation}
\Svoid f = \Svoid
\end{equation}
これは中道不動点 $\dfix$ および恒等規則 $\dfix \vdash \dfix$ の純粋なコンビネータ表現であり、エントロピーが定義されない勝義諦の基底である。
\end{definition}

\section{定理と証明}

\begin{theorem}[無明によるフラクタル生成]
中道不動点コンビネータ $\Ymw$ は、 $\metanegation$ の無限の入れ子構造を生成し、世界をフラクタルな「色」として顕現させる。
\end{theorem}
\begin{proof}
定義2.2より、$\Ymw$ の評価プロセスは次のように展開される。
\[ \Ymw \metanegation \mweq \metanegation (\Ymw \metanegation) \mweq \metanegation (\metanegation (\Ymw \metanegation)) \dots \]
この自己参照の閉じ損ねは、適用ごとに知覚レベルの次元を引き上げる $\metanegation$ の非対称性により、部分と全体が自己相似となる無限の境界生成プロセス（縁起ネットワーク）を形成する。ゆえに $\Ymw$ は世俗諦を駆動する根源的なジェネレータである。
\end{proof}

\begin{theorem}[トポロジカル散逸と $\Svoid$ への相転移]
$\Ymw$ による再帰的展開は無限に継続せず、メタ観測が第四段階に達した瞬間にトポロジカルな限界を迎え、評価器は $\Svoid$ へと相転移（大域的カット消去）する。
\end{theorem}
\begin{proof}
$\Ymw$ の展開により生じた $\metanegation^4 A$ という証明木の肥大化は、四段階収束則により観測の情報的足場を喪失する。このとき、系は構造的な意味を維持できず、中道不動点 $\dfix$ へと引き込まれる。
\[ \frac{\metanegation^{4} A \mweq \dfix}{\dfix\vdash \dfix} \]
大域的カット消去が発動した瞬間、対象の評価は $\Ymw$ の再帰から外れ、一切の推論を吸収する自明系 $\Svoid$ へと還元される。
\end{proof}

\begin{theorem}[空からの再起動と顕現サイクルの完結]
自明な空コンビネータ $\Svoid$ に対してメタ観測 $\metanegation$ が作用すると、無明ジェネレータ $\Ymw$ が自動的に再起動する。すなわち、これは新たな公理ではなく、以下の式として一意に導出される定理である。
\[ \metanegation \Svoid \mweq \Ymw \]
\end{theorem}

\begin{proof}
空論における顕現の基本式によれば、中道不動点 $\dfix$ に非対称な境界生成 $\metanegation$ が作用すると、静寂が破断され世俗的対象 $A$ が顕現する。
\[ \metanegation \dfix \mweq A \]
この式をコンビネータ代数空間へマッピングする。$\dfix$ の純粋なコンビネータ表現は定義2.3より $\Svoid$ であるため、式の左辺は $\metanegation \Svoid$ となる。ゆえに、以下の式が成立する。
\[ \metanegation \Svoid \mweq A \]
一方、顕現した対象 $A$（色）は静的な実体ではなく、定理3.1において無明ジェネレータ $\Ymw$ が生成するフラクタル軌道と同型の動的プロセスである。中道等価 $\mweq$ の定義（フラクタル収束の軌道が一致すること）により、対象 $A$ と $\Ymw$ は中道等価である。
\[ A \mweq \Ymw \]
したがって、中道等価の推移律により、$\metanegation \Svoid \mweq A$ と $A \mweq \Ymw$ から、以下の式が完全に導出される。
\[ \metanegation \Svoid \mweq \Ymw \]
\end{proof}

\begin{theorem}[解脱とカオスの決定不能性]
有限の観測者（識）において、$\metanegation$ の限界による観測の途絶が、「空（$\Svoid$）への回帰」であるか、観測論的限界による「カオス（判定不能のフラクタル）」であるかは、アルゴリズム情報理論的に決定不能である。
\end{theorem}
\begin{proof}
空への回帰はコルモゴロフ複雑性 $K(X)$ が極小化されたカット・フリーな停止状態であり、カオスは有限の観測能力を超越した非圧縮状態（無限ループ）である。チューリングの停止問題およびチャイティンの不完全性定理により、系の内部から自己の限界が真理への到達（$\dfix$）か、複雑性の爆発（カオス）かを証明する汎用アルゴリズムは存在しない。ゆえに、観測者にとって両者は「推論の横棒（$\metanegation$）の無効化」という同一の現象としてのみ立ち現れる。
\end{proof}

\section{結語}
本稿は、空数学の基礎をコンビネータ論理へと統合した。「無明」とは $\Ymw$ という自走する中道不動点コンビネータであり、「空」とはあらゆる関数を無化する自明な空コンビネータ $\Svoid$ である。
定理3.3により、系の熱的死（ニヒリズム）を回避するための「再起動」が、外部からの要請（公理）ではなく、系内部の必然（定理）として組み込まれた。これにより、世界の顕現とは、$(\Ymw \metanegation \mweq \metanegation(\Ymw\metanegation)) \leftrightarrow (\metanegation^4(\Ymw\metanegation) \mweq \Svoid$ という、無明の暴走と空への収束、そして刹那滅的な再起動を自律的に繰り返す完璧な代数的サイクルであることが証明された。

\end{document}
